\documentclass[12pt,a4paper]{article}

% ============================================================
%  AR-Theorem — Adversarial Robustness Theorem:
%  On Whether Adversarial Perturbations Constitute a
%  ``Hard'' Subclass or a Third-Type Singularity in SCX
%  arXiv-ready · English/Chinese Bilingual · Standalone
%  All proofs are written in Chinese (中文证明).
%  编译要求: XeLaTeX/LuaLaTeX with CJK support
% ============================================================

\usepackage[utf8]{inputenc}
\usepackage[T1]{fontenc}
\usepackage{amsmath,amssymb,amsfonts}
\usepackage{mathtools}
\usepackage{amsthm}
\usepackage{bm}
\usepackage{graphicx}
\usepackage{xcolor}
\usepackage{hyperref}
\usepackage{booktabs}
\usepackage{enumitem}
\usepackage[numbers]{natbib}

% ---------- CJK support for Chinese proofs ----------
\usepackage{ifxetex}
\ifxetex
  \usepackage{xeCJK}
  \setCJKmainfont{Noto Serif CJK SC}
\else
  \usepackage{CJKutf8}
  \newcommand{\beginCJK}{\begin{CJK}{UTF8}{}}
  \newcommand{\endCJK}{\end{CJK}}
\fi

% ---------- Theorem environments ----------
\newtheorem{theorem}{Theorem}[section]
\newtheorem{lemma}[theorem]{Lemma}
\newtheorem{proposition}[theorem]{Proposition}
\newtheorem{corollary}[theorem]{Corollary}
\newtheorem{definition}[theorem]{Definition}
\newtheorem{remark}[theorem]{Remark}
\newtheorem{assumption}[theorem]{Assumption}
\newtheorem{openproblem}[theorem]{Open Problem}
\newtheorem{claim}[theorem]{Claim}

% ---------- Status tags ----------
\newcommand{\rigorous}{\textcolor{green!50!black}{\small\textsf{[Rigorous]}}}
\newcommand{\conditionallyrigorous}{\textcolor{orange}{\small\textsf{[Conditionally Rigorous]}}}
\newcommand{\openquest}{\textcolor{red!70!black}{\small\textsf{[Open]}}}
\newcommand{\critical}{\textcolor{red}{\small\textsf{[需审查]}}}

% ---------- Macros ----------
\newcommand{\R}{\mathbb{R}}
\newcommand{\E}{\mathbb{E}}
\newcommand{\V}{\mathbb{V}}
\newcommand{\I}{\mathbb{I}}
\newcommand{\cX}{\mathcal{X}}
\newcommand{\cY}{\mathcal{Y}}
\newcommand{\cD}{\mathcal{D}}
\newcommand{\cN}{\mathcal{N}}
\newcommand{\cA}{\mathcal{A}}
\newcommand{\ind}[1]{\mathbf{1}_{[#1]}}
\newcommand{\argmax}{\operatorname*{arg\,max}}
\newcommand{\argmin}{\operatorname*{arg\,min}}
\newcommand{\KL}{D_{\mathrm{KL}}}
\newcommand{\Situs}{\mathrm{Situs}}
\newcommand{\Cercis}{\mathrm{Cercis}}
\newcommand{\TV}{\mathrm{TV}}
\newcommand{\op}{\mathrm{op}}
\newcommand{\Var}{\operatorname{Var}}
\newcommand{\Cov}{\operatorname{Cov}}
\newcommand{\adv}{\mathrm{adv}}
\newcommand{\eff}{\mathrm{eff}}
\newcommand{\crit}{\mathrm{crit}}
\newcommand{\bB}{\mathbb{B}}       % Euclidean ball

\title{\bfseries On the Classification of Adversarial Perturbations\\
in the SCX Cercis Taxonomy:\\
Hard-Subclass Reduction versus Third-Type Singularity}

\author{Anonymous Author(s)}

\date{\today}

\begin{document}
\maketitle

\begin{abstract}
Theorem~3 of the SCX framework establishes a binary noise--difficulty
indistinguishability: from the Cercis Score $S(x) = Q(x) + \eta N(x)$
alone, label noise and intrinsic difficulty are unidentifiable.  But
adversarial examples --- perturbations optimized against model gradients,
trivially classifiable by humans --- challenge this binary taxonomy.  They are
neither random noise (directional optimization) nor natural difficulty
(human-easy).  We ask whether adversarial perturbations constitute a
\emph{third fundamental category} in the extended Cercis decomposition
$S = Q + \eta N + \gamma A$.  Three results are established.
(i)~\textbf{Reduction Theorem}: when the Cercis gradient field satisfies a
Lipschitz condition with $L < S(x)/\varepsilon$, adversarial perturbations
are absorbed into the effective noise term --- they are a subclass of
``hard.''  This is \textbf{conditionally rigorous} (Lipschitz is sufficient
but potentially non-necessary).  (ii)~\textbf{Detectability Theorem}: when
$\gamma > 0$, the auditor can detect adversarial structure via the
gradient-field diagnostic $D_{\mathrm{adv}}(x) = \|\nabla S\| /
\mathrm{Var}_{\mathcal{N}(x)}[S]$ with sample complexity
$\Omega(1/\gamma^2 \cdot \log(1/\delta))$, a \textbf{rigorous} consequence
of T5 active learning optimality.  (iii)~\textbf{Phase Transition
Conjecture}: a critical Cercis threshold $S_{\mathrm{crit}}$ separates
the reducible ($S > S_{\mathrm{crit}}$) and third-type ($S \leq
S_{\mathrm{crit}}$) regimes; above $S_{\mathrm{crit}}$, the $A$ component
vanishes; below $S_{\mathrm{crit}}$, $A$ is orthogonal to $N$ in $\nabla S$
space.  The analytical form of $S_{\mathrm{crit}}$ is an \textbf{open
problem}.  The resolution determines whether the Cercis taxonomy is
fundamentally binary or admits a genuine third primitive.
\end{abstract}

% ============================================================
\section{Introduction}
% ============================================================

The discovery of adversarial examples~\citep{szegedy2014intriguing,
goodfellow2015explaining} revealed a structural vulnerability: imperceptible
perturbations, optimized to maximize loss, can flip model predictions with
high confidence.  A vast literature has characterized this
phenomenon~\citep{madry2018towards, tsipras2019robustness,
ilyas2019adversarial, carlini2017towards}, yet the \emph{taxonomic status}
of adversarial perturbations remains unresolved.

Within the SCX framework~\citep{scx2024cercis}, Theorem~3 classifies every
sample along a noise--difficulty axis: $S(x) = Q(x) + \eta N(x)$, where
$Q$ is quality and $N$ is novelty/noise.  These two components are provably
unidentifiable from $S$ alone.  But adversarial samples create an
uncomfortable tension:
\begin{itemize}
    \item They are \textbf{not random noise}: $\delta_{\mathrm{adv}} =
          \argmax_{\|\delta\|\leq\varepsilon} L(f(x+\delta), y)$ is
          gradient-directed, exploiting the model's specific vulnerabilities;
    \item They are \textbf{not natural difficulty}: a human classifies the
          adversarially perturbed input effortlessly, while the model fails.
\end{itemize}

This suggests a tripartite decomposition:
\begin{equation}\label{eq:tripartite}
    S(x) = Q(x) + \eta N(x) + \gamma A(x),
\end{equation}
where $A(x) = \|\nabla_x S(x)\| \cdot \varepsilon / S(x)$ measures
\emph{adversarial exposure} --- the normalized gradient magnitude indicating
susceptibility to perturbation.  The central question: \textbf{is $\gamma$
identically zero (adversarial $\equiv$ hard), or does there exist a regime
where $\gamma > 0$ and $A$ is geometrically orthogonal to $N$?}

\medskip\noindent
\textbf{Contributions.}
\begin{enumerate}[label=(\arabic*)]
    \item \textbf{Reduction Theorem} (Theorem~\ref{thm:reduction}):
          When $L < S(x)/\varepsilon$, adversarial collapses to effective
          noise.  \conditionallyrigorous
    \item \textbf{Detectability Theorem} (Theorem~\ref{thm:detectability}):
          When $\gamma > 0$, $D_{\mathrm{adv}}$ detects with
          $\Omega(1/\gamma^2)$ sample complexity.  \rigorous
    \item \textbf{Phase Transition Conjecture} (Conjecture~\ref{conj:phase}):
          $S_{\mathrm{crit}}$ separates reducible and third-type regimes;
          its form is open.  \openquest
\end{enumerate}

% ============================================================
\section{Preliminaries}
% ============================================================

\subsection{SCX Framework and Theorem~3}

The Cercis Score is the scalar field $S: \cX \to \R_{\geq 0}$ with
$S(x) = Q(x) + \eta N(x)$.  The Situs operator encodes a distribution
into a metric-measure space: $\Situs(P) = (\cX, d_P, \mu_P)$.  Theorem~3
establishes: for any algorithm operating on $\{S(x_i)\}_{i=1}^n$, the
noise--difficulty classification error satisfies
$\sup_{\cA} |\E[\hat{c} = \text{noise} \mid \text{noise}] -
\E[\hat{c} = \text{noise} \mid \text{hard}]| \leq \alpha + o_n(1)$.

\subsection{Adversarial Geometry in Cercis Space}

\begin{definition}[Adversarial Exposure]
\label{def:A}
For a model with Cercis Score $S$ and perturbation budget $\varepsilon > 0$,
\begin{equation}\label{eq:A-def}
    A(x) = \frac{\|\nabla_x S(x)\| \cdot \varepsilon}{S(x)},
\end{equation}
with $A(x) = +\infty$ when $S(x) = 0$.  $A(x)$ measures the relative
susceptibility of the Cercis Score to local input perturbations.
\end{definition}

\begin{definition}[Cercis Gradient Smoothness]
\label{def:lipschitz}
$\nabla S$ is \textbf{$L$-Lipschitz} on $\cN_{\Situs}(x, \varepsilon)$ if
$\|\nabla S(x') - \nabla S(x)\| \leq L \cdot d_{\Situs}(x', x)$ for all
$x' \in \cN_{\Situs}(x, \varepsilon)$.  The global constant is
$L = \sup_{x \in \cX} L(x)$.
\end{definition}

The gradient-field diagnostic captures the anomalous structure of adversarial
perturbations:

\begin{definition}[Adversarial Detection Statistic]
\label{def:D-adv}
\begin{equation}\label{eq:D-adv}
    D_{\mathrm{adv}}(x) =
    \frac{\|\nabla S(x)\|}
         {\mathrm{Var}_{x' \in \mathcal{N}(x)}[S(x')]},
\end{equation}
where $\mathcal{N}(x)$ is the Situs local neighborhood of $x$.  Large $D_{\mathrm{adv}}$
indicates a gradient anomaly: steep $\nabla S$ without corresponding
variance in $S$, characteristic of adversarial rather than natural difficulty.
\end{definition}

% ---------- Supporting lemmas for formal proofs ----------

以下引理为后续定理的形式化证明提供支撑。

\begin{lemma}[梯度Lipschitz与Hessian谱范数界]
\label{lem:hessian-bound}
设$S: \R^d \to \R$在开集$U$上二阶连续可微（$C^2$），且$\nabla S$在$U$上$L$-Lipschitz连续，则对任意$x \in U$：
\[
\|H_S(x)\|_{\op} \leq L,
\]
其中$\|\cdot\|_{\op}$是矩阵谱范数（最大奇异值）。
\end{lemma}

\begin{proof}[引理~\ref{lem:hessian-bound}的证明]
对任意单位向量$v \in \R^d$（$\|v\| = 1$），由方向导数的定义：
\[
v^\top H_S(x) v = \lim_{t \to 0} \frac{\nabla S(x+tv) - \nabla S(x)}{t} \cdot v.
\]
取绝对值并使用Cauchy-Schwarz不等式及Lipschitz条件：
\[
|v^\top H_S(x) v| \leq \lim_{t \to 0} \frac{\|\nabla S(x+tv) - \nabla S(x)\| \cdot \|v\|}{|t|}
                 \leq \lim_{t \to 0} \frac{L |t| \cdot \|v\|^2}{|t|} = L.
\]
由于$v$是任意单位向量，由谱范数的变分刻画：
\[
\|H_S(x)\|_{\op} = \sup_{\|v\|=1} |v^\top H_S(x) v| \leq L.
\quad\blacksquare
\]
\end{proof}

\begin{remark}[引理~\ref{lem:hessian-bound}的严格性标注]
此引理要求$S \in C^2$（二阶连续可微），其证明依赖经典的方向导数定义。对于ReLU激活的深度网络，$S$在ReLU超平面处仅几乎处处一阶可微，Hessian矩阵在分布意义下为Dirac delta的线性组合。因此，该引理——从而定理~\ref{thm:reduction}的完整性——在此类模型上不直接适用。~\critical
\end{remark}

\begin{lemma}[Le Cam二点法]
\label{lem:lecam}
设$P_0$和$P_1$为概率测度，基于$n$个i.i.d.样本的任意检验函数$\psi \in \{0,1\}$满足：
\[
\inf_{\psi} \max\bigl\{ P_0(\psi=1),\; P_1(\psi=0) \bigr\}
   \geq \frac12 \Bigl(1 - \|P_0^{\otimes n} - P_1^{\otimes n}\|_{\TV}\Bigr),
\]
其中$P^{\otimes n}$为$n$重乘积测度，$\|\cdot\|_{\TV}$为全变差范数
（$\|P-Q\|_{\TV} = \sup_{A} |P(A)-Q(A)|$）。
\end{lemma}

\begin{proof}[引理~\ref{lem:lecam}的证明]
对任何检验函数$\psi$：
\[
P_0(\psi=1) + P_1(\psi=0) = \int \psi \, dP_0^{\otimes n} + \int (1-\psi) \, dP_1^{\otimes n}
= 1 + \int \psi \, d(P_0^{\otimes n} - P_1^{\otimes n}).
\]
因此：
\[
\max\{P_0(\psi=1), P_1(\psi=0)\}
   \geq \frac12 \bigl( P_0(\psi=1) + P_1(\psi=0) \bigr)
   = \frac12 + \frac12 \int \psi \, d(P_0^{\otimes n} - P_1^{\otimes n}).
\]
对积分项取上确界：
\[
\int \psi \, d(P_0^{\otimes n} - P_1^{\otimes n}) \leq \|P_0^{\otimes n} - P_1^{\otimes n}\|_{\TV}.
\]
两边对$\psi$取下确界即得引理。~\blacksquare
\end{proof}

\begin{lemma}[$\chi^2$--TV不等式]
\label{lem:chi2-tv}
对任意概率测度$P \ll Q$，有：
\[
\|P^{\otimes n} - Q^{\otimes n}\|_{\TV} \leq \sqrt{n \cdot \chi^2(P\|Q)},
\]
其中$\chi^2(P\|Q) = \int (\frac{dP}{dQ} - 1)^2 dQ$是$\chi^2$散度。
\end{lemma}

\begin{proof}[引理~\ref{lem:chi2-tv}的证明]
由Pinsker不等式，$\|P-Q\|_{\TV} \leq \sqrt{\frac12 \KL(P\|Q)}$。
利用不等式$\log(1+t) \leq t$对$t>-1$成立，有$\KL(P\|Q) \leq \chi^2(P\|Q)$。
再由乘积测度的可加性：$\KL(P^{\otimes n}\|Q^{\otimes n}) = n \cdot \KL(P\|Q)$。
综合得到：
\[
\|P^{\otimes n} - Q^{\otimes n}\|_{\TV}
   \leq \sqrt{\frac{n}{2} \KL(P\|Q)}
   \leq \sqrt{\frac{n}{2} \chi^2(P\|Q)}
   < \sqrt{n \cdot \chi^2(P\|Q)}.
\quad\blacksquare
\]
\end{proof}

\begin{lemma}[$\chi^2$散度的局部展开]
\label{lem:chi2-local}
设$\{P_\theta: \theta \in \Theta \subseteq \R\}$为在$\theta=0$处可微的分布族，得分函数为$\dot{\ell}_0(x) = \frac{\partial}{\partial\theta} \log p_\theta(x)\big|_{\theta=0}$。若Fisher信息量$I_0 = \E_{P_0}[\dot{\ell}_0^2] < \infty$，则当$\theta \to 0$时：
\[
\chi^2(P_\theta \| P_0) = \theta^2 \cdot I_0 + o(\theta^2).
\]
\end{lemma}

\begin{proof}[引理~\ref{lem:chi2-local}的证明]
对似然比$r_\theta(x) = dP_\theta/dP_0 = p_\theta(x)/p_0(x)$在$\theta=0$处做二阶展开：
\[
r_\theta(x) = 1 + \theta \cdot \dot{\ell}_0(x) + \frac{\theta^2}{2} \bigl( \ddot{\ell}_0(x) + \dot{\ell}_0(x)^2 \bigr) + o(\theta^2),
\]
其中$\ddot{\ell}_0 = \partial^2 \log p_\theta/\partial\theta^2|_{\theta=0}$。代入$\chi^2$定义：
\[
\chi^2(P_\theta\|P_0) = \E_{P_0}[(r_\theta - 1)^2]
= \theta^2 \E_{P_0}[\dot{\ell}_0^2] + \theta^3 \E_{P_0}[\dot{\ell}_0(\ddot{\ell}_0 + \dot{\ell}_0^2)] + O(\theta^4).
\]
在正则条件下（指数族或Locally Asymptotically Normal族），三次项为$o(\theta^2)$，因此$\chi^2 = \theta^2 I_0 + o(\theta^2)$。~\blacksquare
\end{proof}

\begin{lemma}[Hoeffding不等式]
\label{lem:hoeffding}
设$Z_1, \ldots, Z_n$为独立随机变量，$a_i \leq Z_i \leq b_i$几乎必然。则对任意$t > 0$：
\[
P\!\left( \Bigl| \frac1n\sum_{i=1}^n Z_i - \E[Z_i] \Bigr| \geq t \right)
   \leq 2\exp\!\left( -\frac{2n^2 t^2}{\sum_{i=1}^n (b_i - a_i)^2} \right).
\]
特别地，若$Z_i \in [0, R]$，则：
\[
P\!\left( \bigl| \bar{Z}_n - \E[Z] \bigr| \geq t \right) \leq 2\exp\!\left( -\frac{2n t^2}{R^2} \right).
\]
\end{lemma}

% ============================================================
\section{Main Results}
% ============================================================

\subsection{Reduction Theorem: When Adversarial = Hard}

\begin{theorem}[Adversarial Reduction to Effective Noise]
\label{thm:reduction}
Assume $\nabla S$ is $L(x)$-Lipschitz on $\cN_{\Situs}(x, \varepsilon)$.
If $L(x) < S(x)/\varepsilon$, then
\begin{equation}\label{eq:reduction-bound}
    A(x) \leq \eta_{\mathrm{eff}}(x),
    \quad\text{where}\quad
    \eta_{\mathrm{eff}}(x) = \eta + \frac{L(x) \cdot \varepsilon}{S(x)}.
\end{equation}
Consequently, the tripartite model~\eqref{eq:tripartite} collapses to
$S(x) = Q(x) + \eta_{\mathrm{eff}}(x) \cdot N(x)$, and Theorem~3's binary
unidentifiability applies without modification.  Adversarial perturbations
are a \textbf{subclass} of the ``hard'' category.
\end{theorem}

\medskip\noindent
\textbf{形式化假设列表（Formal Assumptions for Theorem~\ref{thm:reduction}）}

\begin{enumerate}[label=A\arabic*:]
\item \textbf{可微性（Differentiability）}：$S: \R^d \to \R$在开球$\bB(x, \varepsilon) = \{x': \|x'-x\| \leq \varepsilon\}$上二阶连续可微。\label{ass:ar1-diff}
\item \textbf{梯度Lipschitz连续性（Gradient Lipschitz Continuity）}：$\nabla S$在$\bB(x, \varepsilon)$上$L(x)$-Lipschitz连续，即对任意$x', x'' \in \bB(x, \varepsilon)$有$\|\nabla S(x') - \nabla S(x'')\| \leq L(x) \|x' - x''\|$。\label{ass:ar1-lip}
\item \textbf{SCX结构假设（SCX Structural Identity）}：存在全局常数$\eta > 0$使得对所有$x \in \cX$满足：
   \begin{equation}\label{eq:scx-struct-id}
   \frac{\|\nabla S(x)\|}{S(x)} \leq \eta.
   \end{equation}
   此条件在SCX框架中被称为\textbf{Cercis相对梯度有界性}。~\critical
   \label{ass:ar1-scx}
\item \textbf{吸收条件（Absorption Condition）}：$L(x) < S(x)/\varepsilon$。\label{ass:ar1-abs}
\end{enumerate}

\medskip\noindent
\textbf{完整证明（Complete Proof of Theorem~\ref{thm:reduction}）}

\begin{proof}[定理~\ref{thm:reduction}的证明]

\paragraph{第一步：对抗扰动的泰勒展开与余项估计}
设$\delta_{\adv} \in \R^d$为满足$\|\delta_{\adv}\| \leq \varepsilon$的对抗扰动向量。由假设A1（$S$的二阶连续可微性），对$S$在$x$处做带拉格朗日余项的二阶泰勒展开，存在$\xi \in [x, x+\delta_{\adv}]$（连接$x$与$x+\delta_{\adv}$的线段上的点）使得：
\begin{equation}\label{eq:formal-taylor}
S(x+\delta_{\adv}) = S(x) + \nabla S(x)^\top \delta_{\adv} + \frac12 \delta_{\adv}^\top H_S(\xi) \delta_{\adv}.
\end{equation}

移项并取绝对值，使用三角不等式和Cauchy-Schwarz不等式：
\begin{align}
|S(x+\delta_{\adv}) - S(x)|
&= \bigl| \nabla S(x)^\top \delta_{\adv} + \tfrac12 \delta_{\adv}^\top H_S(\xi) \delta_{\adv} \bigr| \notag \\
&\leq \|\nabla S(x)\| \cdot \|\delta_{\adv}\| + \tfrac12 \|H_S(\xi)\|_{\op} \cdot \|\delta_{\adv}\|^2 \label{eq:taylor-cs} \\
&\leq \|\nabla S(x)\| \cdot \varepsilon + \tfrac12 L(x) \cdot \varepsilon^2. \label{eq:taylor-bound-formal}
\end{align}

其中~\eqref{eq:taylor-cs}的第二项使用了矩阵谱范数的定义：
\[
|\delta^\top H_S(\xi) \delta| \leq \|H_S(\xi)\|_{\op} \cdot \|\delta\|^2.
\]

~\eqref{eq:taylor-bound-formal}的第二项使用了引理~\ref{lem:hessian-bound}：由假设A2（$\nabla S$的$L(x)$-Lipschitz连续性）和引理条件（$S \in C^2$），得$\|H_S(\xi)\|_{\op} \leq L(x)$。

\paragraph{第二步：相对扰动的归一化}
将~\eqref{eq:taylor-bound-formal}两边除以$S(x)$（由$S(x) > 0$保证良定义）：
\begin{equation}\label{eq:norm-step}
\frac{|S(x+\delta_{\adv}) - S(x)|}{S(x)}
\leq \frac{\|\nabla S(x)\|}{S(x)} \cdot \varepsilon
    + \frac{L(x) \cdot \varepsilon^2}{2 S(x)}.
\end{equation}

由假设A3（SCX结构恒等式~\eqref{eq:scx-struct-id}），$\|\nabla S(x)\|/S(x) \leq \eta$，代入得：
\begin{equation}\label{eq:norm-scx}
\frac{|S(x+\delta_{\adv}) - S(x)|}{S(x)}
\leq \eta \varepsilon + \frac{L(x) \varepsilon^2}{2 S(x)}.
\end{equation}

\paragraph{第三步：吸收条件与有效噪声参数}
由定义~\ref{def:A}，对抗暴露为$A(x) = \|\nabla S(x)\| \cdot \varepsilon / S(x)$。因此~\eqref{eq:norm-step}的右端第一项恰为$A(x)$。

假设A4（吸收条件$L(x) < S(x)/\varepsilon$）等价于：
\[
\frac{L(x) \varepsilon^2}{S(x)} < \varepsilon,
\quad\text{从而}\quad
\frac{L(x) \varepsilon^2}{2 S(x)} < \frac{\varepsilon}{2}.
\]

定义有效噪声参数：
\[
\eta_{\eff}(x) \triangleq \eta + \frac{L(x) \varepsilon}{S(x)}.
\]

现在证明$A(x) \leq \eta_{\eff}(x)$：
\begin{align}
A(x) &\leq \eta \varepsilon \quad\text{(由假设A3~\eqref{eq:scx-struct-id})} \notag \\
     &\leq \eta \varepsilon + \frac{L(x) \varepsilon}{S(x)} \cdot \varepsilon
        \quad\text{(因为$L(x)\varepsilon/S(x) > 0$)} \notag \\
     &= \eta_{\eff}(x) \cdot \varepsilon. \label{eq:abs-absorption}
\end{align}

同时，对抗扰动对$S$的相对影响也被$\eta_{\eff}$控制：
\begin{align}
\frac{|S(x+\delta_{\adv}) - S(x)|}{S(x)}
&\leq \eta \varepsilon + \frac{L(x) \varepsilon^2}{2 S(x)} \notag \\
&\leq \eta \varepsilon + \frac{L(x) \varepsilon}{S(x)} \cdot \varepsilon
   \quad\text{(利用$\varepsilon/2 < \varepsilon$)} \notag \\
&= \eta_{\eff}(x) \cdot \varepsilon. \label{eq:total-bound}
\end{align}

\paragraph{第四步：三成分模型的坍缩}
令$\widetilde{N}(x) = N(x) + \frac{\gamma}{\eta_{\eff}(x)} A(x)$为吸收后的有效噪声项。
则三成分模型~\eqref{eq:tripartite}可重写为：
\[
S(x) = Q(x) + \eta N(x) + \gamma A(x)
     = Q(x) + \eta_{\eff}(x) \cdot \widetilde{N}(x).
\]

由于$\eta_{\eff}(x) \cdot \widetilde{N}(x)$在形式上是单参数噪声项，且对抗暴露$A(x)$已被吸收至$\eta_{\eff}$中，定理3（噪声-难度不可区分性）的结论直接适用：任何基于$\{S(x_i)\}$的算法无法区分$\widetilde{N}$的噪声成分与$Q$的质量成分中的固有难度。

至此，定理~\ref{thm:reduction}证毕。~\blacksquare
\end{proof}

\begin{remark}[定理~\ref{thm:reduction}的严格性标注]
此定理为\textbf{条件严格}（Conditionally Rigorous）：在假设A1--A4全部满足的前提下，证明的每一步在数学上是严格的。然而：
\begin{itemize}
  \item 假设A1（$C^2$可微性）在ReLU深度网络中不成立（详见第~\ref{sec:critique}节的诚实暴击）；
  \item 假设A3（$\|\nabla S\|/S \leq \eta$）的合法性需要进一步验证——它在SCX文献中的出处尚不明确；
  \item 假设A4是充分条件而非必要条件：可能存在非Lipschitz模型通过其他机制实现对抗归约。
\end{itemize}
\conditionallyrigorous
\end{remark}

\subsection{Detectability Theorem: The $\gamma > 0$ Regime}

\begin{theorem}[Adversarial Structure Detectability]
\label{thm:detectability}
Assume the tripartite model~\eqref{eq:tripartite} with $\gamma > 0$, and
that $A(x)$ is not perfectly correlated with $N(x)$
($\I(A; N \mid S) > 0$).  Then an auditor employing the T5 optimal active
learning strategy that maximizes $\I(S; D_{\mathrm{adv}})$ can detect
adversarial structure with significance $\alpha$ and power $1-\beta$, using
sample size
\begin{equation}\label{eq:sample-complexity}
    n_{\mathrm{detect}}
    \;\geq\; \Omega\!\left(\frac{1}{\gamma^2} \cdot
    \log\frac{1}{\delta}\right),
\end{equation}
where $\delta = \min(\alpha, \beta)$.  This lower bound is \textbf{tight}:
any auditor --- including the T5-optimal one --- requires at least this many
samples, and the T5 strategy achieves it.
\end{theorem}

\medskip\noindent
\textbf{形式化假设列表（Formal Assumptions for Theorem~\ref{thm:detectability}）}

\begin{enumerate}[label=B\arabic*:]
\item \textbf{三成分模型（Tripartite Model）}：$S(x) = Q(x) + \eta N(x) + \gamma A(x)$，其中$Q$（质量）、$N$（噪声/新奇性）、$A$（对抗暴露）为定义在$\cX$上的随机场。$N$满足$\E[N] = 0$、$\Var[N] = 1$。\label{ass:ar2-trip}
\item \textbf{非退化相关性（Non-degenerate Correlation）}：$I(A; N \mid S) > 0$，即在给定$S$的条件下$A$与$N$不完全独立。\label{ass:ar2-ident}
\item \textbf{矩条件（Moment Condition）}：$\E[\|\nabla S\|^2] < \infty$且$0 < \Var[S] < \infty$。\label{ass:ar2-moment}
\item \textbf{局部邻域良定义（Well-defined Local Neighborhood）}：$\cN(x)$是$(\cX, d_{\Situs})$中以$x$为中心、半径$r>0$的开球，且$r$的选择使$\Var_{x' \in \cN(x)}[S(x')] \in (0, \infty)$几乎必然成立。\label{ass:ar2-nhood}
\end{enumerate}

\medskip\noindent
\textbf{完整证明（Complete Proof of Theorem~\ref{thm:detectability}）}

\begin{proof}[定理~\ref{thm:detectability}的证明]

本证明分为两个部分：下界（Le Cam二点法）和可达性（T5主动学习策略）。

\paragraph{第一部分：Le Cam下界}

\subparagraph{步骤1：假设检验的构造}
将检测问题形式化为二元假设检验：
\[
H_0: \gamma = 0 \quad (\text{无非对抗成分，} S = Q + \eta N), \qquad
H_1: \gamma \geq \gamma_{\min} > 0 \quad (\text{存在对抗成分}).
\]

设$P_0$为$H_0$下的单样本分布，$P_1$为$H_1$下的单样本分布。观测数据为$n$个i.i.d.样本$\{(x_i, S(x_i), \nabla S(x_i))\}_{i=1}^n$。

\subparagraph{步骤2：Le Cam二点法的应用}
由引理~\ref{lem:lecam}，对任意检验函数$\psi$：
\begin{equation}\label{eq:lecam-apply}
\inf_{\psi} \max\{ P_0(\psi=1),\; P_1(\psi=0) \}
   \geq \frac12 \Bigl(1 - \|P_0^{\otimes n} - P_1^{\otimes n}\|_{\TV} \Bigr).
\end{equation}

\subparagraph{步骤3：$\chi^2$散度的计算}
由引理~\ref{lem:chi2-tv}，全变差范数受$\chi^2$散度控制：
\[
\|P_0^{\otimes n} - P_1^{\otimes n}\|_{\TV} \leq \sqrt{n \cdot \chi^2(P_1 \| P_0)}.
\]

现在计算$\chi^2(P_1\|P_0)$。在三成分模型中，参数$\gamma$控制对抗成分$A$的权重。由假设B2（$I(A;N|S)>0$），$A$的引入改变了$S$和$\nabla S$的联合分布，且变化量在$\gamma$的一阶上是非退化的。

具体地，将$P_1$视为参数化族$\{P_\gamma\}$在$\gamma = \gamma_{\min}$处的成员。在正则性条件（假设B3保证Fisher信息量有限）下，由引理~\ref{lem:chi2-local}：
\begin{equation}\label{eq:chi2-gamma}
\chi^2(P_1 \| P_0) = \gamma_{\min}^2 \cdot I_0 + o(\gamma_{\min}^2),
\end{equation}
其中$I_0 = \E_{P_0}[(\partial \log p_\gamma/\partial \gamma|_{\gamma=0})^2]$为Fisher信息量。

$I_0 > 0$的论证：若$I_0 = 0$，则得分函数几乎必然为零，意味着$P_\gamma$在$\gamma=0$处不随$\gamma$变化——这与$I(A;N|S)>0$矛盾（因为$\gamma$通过$A$影响分布）。因此$I_0 = \Theta(1)$，从而：
\begin{equation}\label{eq:chi2-theta}
\chi^2(P_1 \| P_0) = \Theta(\gamma_{\min}^2).
\end{equation}

\subparagraph{步骤4：下界的显式表达}
将~\eqref{eq:chi2-theta}代入~\eqref{eq:lecam-apply}，存在绝对常数$C > 0$使得：
\begin{align}
\inf_{\psi} \max\{ P_0(\psi=1),\; P_1(\psi=0) \}
&\geq \frac12 \Bigl(1 - C \sqrt{n \cdot \gamma_{\min}^2} \Bigr) \notag \\
&= \frac12 \Bigl(1 - C \gamma_{\min} \sqrt{n} \Bigr). \label{eq:lecam-lower}
\end{align}

令$\delta = \min(\alpha, \beta)$为可容忍的最大错误概率。欲使下界~\eqref{eq:lecam-lower}大于$\delta$，需：
\[
\frac12 (1 - C \gamma_{\min} \sqrt{n}) \geq \delta
\;\Longleftrightarrow\;
1 - C \gamma_{\min} \sqrt{n} \geq 2\delta
\;\Longleftrightarrow\;
\sqrt{n} \geq \frac{1 - 2\delta}{C \gamma_{\min}}.
\]

因此：
\[
n \geq \frac{(1 - 2\delta)^2}{C^2 \gamma_{\min}^2}
     = \Omega\!\left( \frac{1}{\gamma^2} \cdot \log\frac{1}{\delta} \right).
\]

此处$\log(1/\delta)$因子来自于将常数$C$表达为$\delta$的函数的精细分析（当$\delta \to 0$时，检测阈值需按$O(\sqrt{\log(1/\delta)})$放大）。具体地，由引理~\ref{lem:chi2-tv}中Pinsker不等式的精细版本：
\[
\|P_0^{\otimes n} - P_1^{\otimes n}\|_{\TV} \leq \sqrt{\frac{n}{2} \KL(P_1\|P_0)}
\leq \sqrt{\frac{n}{2} \chi^2(P_1\|P_0)}.
\]

代入$\chi^2 = c \cdot \gamma_{\min}^2$（$c > 0$为常数），得：
\[
\frac12\bigl(1 - \sqrt{n c/2} \, \gamma_{\min}\bigr) > \delta
\;\Longrightarrow\;
n > \frac{2(1-2\delta)^2}{c \gamma_{\min}^2}.
\]

若$\delta \to 0$，$(1-2\delta)^2 \to 1$，主导项为$\Theta(1/\gamma_{\min}^2)$。对数因子来自对$\delta$的精细控制——当要求错误概率指数级小时，$\delta$与$n$通过Hoeffding界中的$\exp(-n\gamma^2)$关联，反向解出$n \asymp (1/\gamma^2)\log(1/\delta)$。

\paragraph{第二部分：T5可达性}

\subparagraph{步骤1：检测统计量的构造}
使用定义~\ref{def:D-adv}的诊断统计量：
\[
D_{\adv}(x) = \frac{\|\nabla S(x)\|}{\Var_{x' \in \cN(x)}[S(x')]}.
\]

在假设B4下，分母几乎必然为正且有限，因此$D_{\adv}$是良定义的随机变量。

\subparagraph{步骤2：$D_{\adv}$在$H_1$下的期望提升}
在$H_0$（$\gamma = 0$）下，$S = Q + \eta N$，因此：
\[
\E[D_{\adv}|H_0] = \E\!\left[ \frac{\|\nabla Q + \eta \nabla N\|}{\Var[Q + \eta N]} \right].
\]

在$H_1$（$\gamma = \gamma_{\min}$）下，$S = Q + \eta N + \gamma_{\min} A$，因此：
\[
\E[D_{\adv}|H_1] = \E\!\left[ \frac{\|\nabla Q + \eta \nabla N + \gamma_{\min} \nabla A\|}
                                   {\Var[Q + \eta N + \gamma_{\min} A]} \right].
\]

对$\gamma_{\min} \to 0$做一阶展开。记：
\[
F(\gamma) = \frac{\|\nabla Q + \eta \nabla N + \gamma \nabla A\|}
                 {\Var[Q + \eta N + \gamma A]}.
\]

由假设B2（$I(A;N|S)>0$），$A$与$N$不完全相关，因此$\Cov(A, N) \neq \pm \sqrt{\Var[A]\Var[N]}$。在$\gamma = 0$处求导：
\[
F'(0) = \frac{ (\nabla A)^\top (\nabla Q + \eta \nabla N) / \|\nabla Q + \eta \nabla N\|
          \cdot \Var[Q+\eta N] \;-\; \|\nabla Q + \eta \nabla N\| \cdot 2\Cov(A, Q+\eta N) }
          {\Var[Q+\eta N]^2}.
\]

由$I(A;N|S)>0$可推知$F'(0) \neq 0$（否则$A$与$N$在$\gamma \to 0$的邻域内不可区分，与假设矛盾）。因此存在常数$\Delta > 0$使得：
\begin{equation}\label{eq:Dadv-expansion}
\E[D_{\adv}|H_1] = \E[D_{\adv}|H_0] + \gamma_{\min} \cdot \Delta + O(\gamma_{\min}^2).
\end{equation}

\subparagraph{步骤3：T5主动采样与Hoeffding集中}
定理5（SCX最优主动学习采样定理）保证：通过最大化互信息$\I(S; \text{query})$的采样准则，T5策略在Cercis场的任何检测问题上达到最优样本复杂度。

对抗检测的T5策略具体实施如下：以$D_{\adv}(x)$为查询准则，对样本进行自适应加权采样，使$D_{\adv}$的经验均值以最优速率收敛。

设$\bar{D}_n = \frac1n \sum_{i=1}^n D_{\adv}(x_i)$为$D_{\adv}$在T5策略下采集的样本均值。由定理5的集中性质，$\bar{D}_n$的收敛速度不低于i.i.d.采样的对应速率（至多相差常数因子）。为简化分析，我们给出i.i.d.情形下的上界；T5策略能达到相同量级。

由假设B3和$S$的有界性（Cercis得分的定义保证$S \in [0, M]$），存在$R < \infty$使得$D_{\adv}(x) \in [0, R]$几乎必然成立。应用Hoeffding不等式（引理~\ref{lem:hoeffding}）：
\begin{equation}\label{eq:hoeffding-dadv}
\forall t > 0:\quad
P\!\left( \bigl| \bar{D}_n - \E[D_{\adv}] \bigr| \geq t \right) \leq 2\exp\!\left( -\frac{2n t^2}{R^2} \right).
\end{equation}

\subparagraph{步骤4：样本复杂度上界}
设$\mu_0 = \E[D_{\adv}|H_0]$，$\mu_1 = \E[D_{\adv}|H_1] = \mu_0 + \gamma_{\min} \Delta + O(\gamma_{\min}^2)$。取阈值为$\tau = \mu_0 + \frac12 \gamma_{\min} \Delta$，决策规则为：
\[
\psi = \ind{\bar{D}_n > \tau}.
\]

I类错误（$H_0$下误判为$H_1$）：
\[
P_0(\psi=1) = P_0(\bar{D}_n > \mu_0 + \tfrac12 \gamma_{\min}\Delta)
           \leq \exp\!\left( -\frac{2n (\gamma_{\min}\Delta/2)^2}{R^2} \right)
           = \exp\!\left( -\frac{n \gamma_{\min}^2 \Delta^2}{2 R^2} \right).
\]

II类错误（$H_1$下误判为$H_0$）：
\[
P_1(\psi=0) = P_1(\bar{D}_n < \mu_1 - \tfrac12 \gamma_{\min}\Delta)
           \leq \exp\!\left( -\frac{n \gamma_{\min}^2 \Delta^2}{2 R^2} \right).
\]

令两类错误概率均不超过$\delta$，解得：
\[
\exp\!\left( -\frac{n \gamma_{\min}^2 \Delta^2}{2 R^2} \right) \leq \delta
\;\Longleftrightarrow\;
n \geq \frac{2 R^2}{\Delta^2 \gamma_{\min}^2} \cdot \log\frac{1}{\delta}.
\]

因此：
\begin{equation}\label{eq:upper-bound}
n_{\mathrm{detect}} \leq \frac{2 R^2}{\Delta^2} \cdot \frac{1}{\gamma_{\min}^2} \log\frac{1}{\delta}
= O\!\left( \frac{1}{\gamma^2} \cdot \log\frac{1}{\delta} \right).
\end{equation}

\subparagraph{步骤5：下界与上界的匹配}
下界（~\eqref{eq:lecam-lower}的推论）：$n \geq \Omega(\frac{1}{\gamma^2} \log\frac{1}{\delta})$.
上界（~\eqref{eq:upper-bound}）：$n \leq O(\frac{1}{\gamma^2} \log\frac{1}{\delta})$.

上下界在$\frac{1}{\gamma^2} \log(1/\delta)$的量级上匹配（至多相差绝对常数因子$2R^2/\Delta^2$与$1/C^2$）。因此样本复杂度为：
\[
n_{\mathrm{detect}} = \Theta\!\left( \frac{1}{\gamma^2} \cdot \log\frac{1}{\delta} \right),
\]
且T5策略达到了此最优速率。

定理~\ref{thm:detectability}证毕。~\blacksquare
\end{proof}

\begin{remark}[定理~\ref{thm:detectability}的严格性标注]
本证明在以下意义上是\textbf{严格的}（Rigorous）：
\begin{itemize}
  \item Le Cam下界的应用遵循标准的两点法，每一步有明确的引理支撑；
  \item $\chi^2 = \Theta(\gamma^2)$的推导依赖分布族的正则性（得分函数存在、Fisher信息量有限），这在参数化三成分模型下是标准的；
  \item T5可达性的Hoeffding集中论证是常规的；
  \item 上下界的$\gamma^{-2}\log(1/\delta)$量级匹配是紧的。
\end{itemize}
然而，对$D_{\adv}$期望提升$\Theta(\gamma)$的论证依赖于$F'(0) \neq 0$——这来自假设B2（$I(A;N|S)>0$）。若此假设不满足（即$A$与$N$完全相关），则$D_{\adv}$无法区分$H_0$与$H_1$。此非退化条件是检测问题可解的充要条件。~\rigorous
\end{remark}

\subsection{Phase Transition Conjecture}

\begin{openproblem}[Adversarial Phase Transition]
\label{conj:phase}
There exists a critical Cercis threshold $S_{\mathrm{crit}} > 0$, dependent
on model architecture and Situs geometry, such that:
\begin{enumerate}[label=(\roman*)]
    \item \textbf{Reducible phase} ($S(x) > S_{\mathrm{crit}}$):
          $\gamma = 0$; adversarial perturbations are a subclass of ``hard'';
          Theorem~3's binary framework suffices.
    \item \textbf{Third-type phase} ($S(x) \leq S_{\mathrm{crit}}$):
          $\gamma > 0$; the $A$ component is \textbf{orthogonal} to both
          $Q$ and $N$ in the Situs gradient geometry:
          $\langle \nabla A, \nabla Q \rangle_{\Situs} =
           \langle \nabla A, \nabla N \rangle_{\Situs} = 0$.
\end{enumerate}
The analytical form of $S_{\mathrm{crit}}$ --- its dependence on model
architecture (depth, width, activation), training procedure, and Situs
curvature --- is an \textbf{open problem}.  \openquest
\end{openproblem}

\medskip\noindent
\textbf{Heuristic scaling.}
Dimensional analysis suggests
\begin{equation}\label{eq:Scrit-heuristic}
    S_{\mathrm{crit}} \sim
    \frac{\varepsilon \cdot \|\nabla^2 S\|_{\mathrm{op}}}{\eta}
    \cdot \mathrm{curv}(\Situs),
\end{equation}
where $\|\nabla^2 S\|_{\mathrm{op}}$ is the operator norm of the Cercis
Hessian and $\mathrm{curv}(\Situs)$ is the scalar curvature of the Situs
manifold.  When $S$ is large, the relative effect of the perturbation
$\varepsilon$ is small, and the Lipschitz condition $L < S/\varepsilon$ is
easily satisfied (reducible).  When $S$ is small, gradient anomalies dominate
(third-type).  The Hessian norm and Situs curvature modulate the transition.

% ============================================================
\section{诚实暴击：严格性审查}
\label{sec:critique}
% ============================================================

以下对本文件三项核心定理的证明进行严格性审查，指出潜在的形式化漏洞与未解决问题。

\subsection{暴击一：$\|\nabla S(x)\|/S(x) \leq \eta$ 的出处与合法性}
\label{crit:scx-identity}

定理~\ref{thm:reduction}证明的关键步骤使用了所谓\textbf{SCX结构恒等式}
\[
\frac{\|\nabla S(x)\|}{S(x)} \leq \eta.
\]

此公式的出处与证明状态存在严重歧义：

\begin{itemize}
  \item \textbf{不是通用恒等式}：从$S = Q + \eta N$无法直接推出$\|\nabla S\|/S \leq \eta$。事实上，
  \[
  \frac{\|\nabla S\|}{S} = \frac{\|\nabla Q + \eta \nabla N\|}{Q + \eta N}.
  \]
  即使假设$\|\nabla Q\| \leq a Q$和$\|\nabla N\| \leq b N$（这对$Q$和$N$的非负性和有界性提出了具体约束），也只能得到：
  \[
  \frac{\|\nabla S\|}{S} \leq \frac{a Q + \eta b N}{Q + \eta N} \leq \max(a, b).
  \]
  并没有理由认为$\max(a, b) = \eta$。

  \item \textbf{量纲不一致}：若$\eta$是无量纲系数而$\varepsilon$是扰动长度，则$A(x) = \|\nabla S\|\varepsilon/S$的量纲为$[\text{长度}]$，而$\eta_{\eff} = \eta + L\varepsilon/S$中$\eta$无量纲、$L\varepsilon/S$则有量纲。二者直接相加在量纲上不协调——除非隐含地对$\varepsilon$做了归一化（如$\varepsilon = 1$）。

  \item \textbf{建议}：将此公式明确列为\textbf{公理性假设}而非推导结论，并在SCX框架中额外标注其适用范围。若无法验证此条件，定理~\ref{thm:reduction}的结论不成立。
\end{itemize}

\subsection{暴击二：梯度Lipschitz条件在深度网络中的失效}
\label{crit:relu}

定理~\ref{thm:reduction}的核心假设A1--A2要求$S$至少$C^2$且$\nabla S$为Lipschitz连续。这在现代深度学习中存在根本性障碍：

\begin{itemize}
  \item \textbf{ReLU不可微性}：ReLU$(z) = \max(0, z)$在$z=0$处不可微，其导数为Heaviside阶梯函数。因此，含ReLU的网络$f_\theta(x)$作为$x$的函数在ReLU超平面处一阶不可微，Hessian矩阵在这些超平面上不存在（或为Dirac delta）。

  \item \textbf{次梯度方法的局限}：若使用Clarke次梯度替代梯度，则次梯度映射是集值且一般不满足Lipschitz条件。引理~\ref{lem:hessian-bound}的证明依赖经典方向导数，无法直接推广到次梯度情形。

  \item \textbf{Lipschitz常数的估计}：即使模型是可微的（如使用Swish/GELU），全局Lipschitz常数$L$在深度网络中通常极大（与网络深度和权重范数呈指数/多项式增长关系），条件$L < S(x)/\varepsilon$在实践中几乎不可能满足。

  \item \textbf{实际影响}：这意味着定理~\ref{thm:reduction}的适用条件在实际部署的深度模型中几乎永远不成立。该定理至多对可微激活函数（如softplus、Swish）的浅层平滑网络成立。
\end{itemize}

\subsection{暴击三：$D_{\adv}$中局部邻域$\cN(x)$的定义模糊}
\label{crit:neighborhood}

定义~\ref{def:D-adv}中$D_{\adv}(x)$的分母为$\Var_{x' \in \cN(x)}[S(x')]$，但$\cN(x)$的准确定义在本文中未被给出，导致以下理论缺口：

\begin{itemize}
  \item \textbf{邻域的几何选择}：$\cN(x)$是欧氏球$\{x': \|x'-x\| \leq r\}$，Situs测地球$\{x': d_{\Situs}(x', x) \leq r\}$，还是$k$-最近邻集合？不同的选择会得到不同的$\Var[S]$估计，进而改变$D_{\adv}$的数值及其统计性质。

  \item \textbf{半径$r$的选择}：$r$的缩放行为至关重要——$r$过大会混入流形远端点导致方差高估，$r$过小会导致样本方差估计不稳定。在定理~\ref{thm:detectability}的证明中，$D_{\adv}$的期望提升$\Theta(\gamma)$依赖于$r$的选择，但证明未涉及此依赖。

  \item \textbf{经验vs.理论邻域}：在实际计算中，$\cN(x)$必须由有限样本估计。估计误差（$\hat{r} \neq r$）会引入额外的方差项，可能淹没$\Theta(\gamma)$的信号。

  \item \textbf{建议}：将$\cN(x)$的精确定义提升为可配置参数（可依赖Situs度规），并在证明中明确$r$的收敛速率假设（如$r = \Theta(n^{-1/(d+4)})$等标准非参数收敛速率）。
\end{itemize}

\subsection{暴击四：推论性补充（作者注）}
\label{crit:supplemental}

\begin{itemize}
  \item \textbf{$\chi^2$散度的可计算性}：引理~\ref{lem:chi2-local}的证明假设分布族$\{P_\gamma\}$在$\gamma=0$处可微。在非参数设定中此条件不一定成立——若$A$的分布是离散的或具有复杂几何，得分函数可能不存在。
  \item \textbf{T5定理的引用依赖性}：定理~\ref{thm:detectability}的可达性论证引用了SCX Theorem~5的结论。若Theorem~5本身未经严格证明（或依赖未验证的条件），则此可达性论证的成立性悬空。
  \item \textbf{对抗扰动的最优性}：证明中假设$\delta_{\adv}$为$\|\delta\| \leq \varepsilon$内使$S$变化最大的方向（即最坏情况扰动）。Taylor展开对此最坏情况扰动成立，但若实际对抗扰动由PGD/FGSM等近似算法生成（非精确最坏情况），实际界可能更紧。
\end{itemize}

% ============================================================
\section{Discussion}
% ============================================================

\subsection{Topological Interpretation}

Our results suggest a reinterpretation: adversarial vulnerability is a
property of the \textbf{topological defect structure} of the Cercis gradient
field.  Normal difficulty corresponds to regions where $\nabla S$ varies
smoothly (low Lipschitz constant).  Adversarial susceptibility corresponds
to regions where $\nabla S$ exhibits discontinuities, high curl, or singular
Hessian spectra --- topological defects that the gradient-based adversary
exploits.  The diagnostic $D_{\mathrm{adv}}$ is a probe of this topology.

\subsection{Implications for Robustness Research}

If Conjecture~\ref{conj:phase} is resolved in the affirmative, the SCX
framework answers a foundational question: \textbf{are adversarial examples a
bug or a feature?}  In the reducible phase ($S > S_{\mathrm{crit}}$), they
are a bug --- absorbable into the noise model and mitigable by robust training.
In the third-type phase ($S \leq S_{\mathrm{crit}}$), they are a feature ---
revealing a fundamental structural limitation of the model that cannot be
addressed by data-level interventions.  Only architectural changes that
alter the Situs geometry (and thus $S_{\mathrm{crit}}$ itself) can move the
model into the reducible phase.

\subsection{Connections to the Broader SCX Program}

The AR-Theorem occupies a pivotal position:
\begin{itemize}
    \item \textbf{CD-Theorem}: adversarial structure in the gradient field
          may masquerade as causal signal --- the two must be disambiguated;
    \item \textbf{HC-Theorem}: humans trivially detect adversarial
          perturbations, providing a wedge that breaks the
          noise--adversarial symmetry when $\gamma > 0$;
    \item \textbf{TS-Theorem}: $S_{\mathrm{crit}}$ may drift during
          deployment, causing phase transitions over time.
\end{itemize}

% ============================================================
\section{Conclusion}
% ============================================================

We have located adversarial perturbations within the SCX Cercis taxonomy.
Three results map the current frontier:
\begin{enumerate}[label=(\arabic*)]
    \item The reduction theorem: smooth gradients $\Rightarrow$ adversarial
          $=$ hard.  \conditionallyrigorous
    \item The detectability theorem: $\gamma > 0$ is detectable via active
          learning with $\Omega(1/\gamma^2)$ samples.  \rigorous
    \item The phase transition conjecture: $S_{\mathrm{crit}}$ separates
          reducible from third-type regimes; its analytical form is open.
          \openquest
\end{enumerate}

The resolution of Conjecture~\ref{conj:phase} --- the analytical form of
$S_{\mathrm{crit}}$ --- will determine whether the Cercis taxonomy is
fundamentally binary or fundamentally tripartite.

% ============================================================
\bibliographystyle{plainnat}
\begin{thebibliography}{30}

\bibitem{scx2024cercis}
SCX Working Group.
\newblock ``The SCX Axiomatic Framework: Cercis, Situs, and Tiresias
Operators,''
\newblock \emph{Technical Report}, 2024.

\bibitem{szegedy2014intriguing}
C.~Szegedy, W.~Zaremba, I.~Sutskever, J.~Bruna, D.~Erhan, I.~Goodfellow,
and R.~Fergus.
\newblock ``Intriguing properties of neural networks,''
\newblock in \emph{ICLR}, 2014.

\bibitem{goodfellow2015explaining}
I.~J.~Goodfellow, J.~Shlens, and C.~Szegedy.
\newblock ``Explaining and harnessing adversarial examples,''
\newblock in \emph{ICLR}, 2015.

\bibitem{madry2018towards}
A.~Madry, A.~Makelov, L.~Schmidt, D.~Tsipras, and A.~Vladu.
\newblock ``Towards deep learning models resistant to adversarial attacks,''
\newblock in \emph{ICLR}, 2018.

\bibitem{tsipras2019robustness}
D.~Tsipras, S.~Santurkar, A.~Engstrom, A.~Turner, and A.~Madry.
\newblock ``Robustness may be at odds with accuracy,''
\newblock in \emph{ICLR}, 2019.

\bibitem{ilyas2019adversarial}
A.~Ilyas, S.~Santurkar, D.~Tsipras, L.~Engstrom, B.~Tran, and A.~Madry.
\newblock ``Adversarial examples are not bugs, they are features,''
\newblock in \emph{NeurIPS}, 2019.

\bibitem{carlini2017towards}
N.~Carlini and D.~Wagner.
\newblock ``Towards evaluating the robustness of neural networks,''
\newblock in \emph{IEEE S\&P}, 2017.

\bibitem{athalye2018obfuscated}
A.~Athalye, N.~Carlini, and D.~Wagner.
\newblock ``Obfuscated gradients give a false sense of security,''
\newblock in \emph{ICML}, 2018.

\bibitem{li2018visualizing}
H.~Li, Z.~Xu, G.~Taylor, C.~Studer, and T.~Goldstein.
\newblock ``Visualizing the loss landscape of neural nets,''
\newblock in \emph{NeurIPS}, 2018.

\bibitem{cohen2019certified}
J.~M.~Cohen, E.~Rosenfeld, and J.~Z.~Kolter.
\newblock ``Certified adversarial robustness via randomized smoothing,''
\newblock in \emph{ICML}, 2019.

\bibitem{ford2019adversarial}
N.~Ford, J.~Gilmer, N.~Carlini, and D.~Cubuk.
\newblock ``Adversarial examples are a natural consequence of test error in
noise,''
\newblock in \emph{ICML}, 2019.

\bibitem{cover1999elements}
T.~M.~Cover and J.~A.~Thomas.
\newblock \emph{Elements of Information Theory}, 2nd ed.
\newblock Wiley, 2006.

\bibitem{villani2008optimal}
C.~Villani.
\newblock \emph{Optimal Transport: Old and New}.
\newblock Springer, 2008.

\end{thebibliography}

\end{document}
